% ============================================================
% One-page progress summary — compiles as-is on Overleaf.
%
% Upload the following image to the Overleaf project:
%   pic2.png — per-tile QP representations
%
% Figures use [H], so they remain exactly where inserted.
% If the document exceeds one page, change 10pt to 9pt.
% ============================================================

\documentclass[10pt,conference]{IEEEtran}

\usepackage{graphicx}
\usepackage{float}      % Provides the [H] placement option
\usepackage{enumitem}
\usepackage{booktabs}
\usepackage{amsmath}
\usepackage{amssymb}


\setlist{
    noitemsep,
    topsep=2pt,
    partopsep=0pt,
    leftmargin=*
}

\begin{document}

\title{Summary and Questions}

\author{\IEEEauthorblockN{Jad Al Lakkis}}


\maketitle
\thispagestyle{empty}

\begin{abstract}This is a quick summary of the progress in data exploration since last meeting and some questions about modeling choices I encountered.
\end{abstract}

\section{Summary}

\begin{itemize}
    \item Extracted the dataset
    (\texttt{hn.mat}, \texttt{rd.mat}, 15 videos) , analzed it , and encountered some missing things that I have to choose and add before implementing the actual code.

\end{itemize}


\section{Key Open Questions and Modeling Decisions}

\noindent
\textbf{Channel-gain parameters.}
\noindent
 This dataset doesnt provide anything about channel gain. But the related 2025 paper provides the following closed-form Rician-fading and path-loss channel model:

\[
h^t=
\frac{3.24\times10^{-4}|\psi^t|^2}
{\left[d^t\tan^{-1}\left(\frac{24\sqrt{2}}{5}\right)\right]^2},
\qquad
K=12~\mathrm{dB},
\qquad
\mathbb{E}[|\psi^t|^2]=1.
\]
\textbf{Should this model be reused as-is?}

\medskip

\noindent
\textbf{Base/enhancement-layer implementation.}
The dataset has no scalable-video-coding layers (no BL-EL split) , only independent QP
representations.
\begin{figure}[H]
    \centering
    \includegraphics[width=0.90\linewidth]{pic2.png}
    \caption{Per-tile QP representations available in the dataset (shown only 8 of them in the screenshot)}
    \label{fig:qp_representations}
\end{figure}

The proposed approach sends a base-quality representation
for every tile and models an enhanced tile by replacing its base-quality
representation with a higher-quality version, rather than transmitting two
complete streams. So the data-volume

\[
A^t
=
A_{\text{base}}^t
+
\sum_i x_i^t\Delta A_i^t,
\qquad
\Delta A_i^t
=
A_{i,\text{selected}}^t-A_{i,\text{base}}^t.
\]

\noindent
\textbf{Does this match the intended formulation, or is genuine scalable
and additive enhancement data required? Which QP indices should represent
the base and enhanced qualities?}

\begin{table}[H]
    \centering
    \footnotesize
    \begin{tabular}{@{}p{2.1cm}p{5.3cm}@{}}
        \toprule
        \textbf{Decision} & \textbf{Options} \\
        \midrule

        FOV source &
        To define the actual Viewport, I need the yaw, pitch and FOV.. do i define the Fov the same as the 2025 paper?\\

        UAV mobility &
        Do i use the same static-distance $d=50\,\mathrm{m}$ case
        reported in the 2025 paper? \\

        Training &
        Do I train on all 15 videos, or choose one of them at first and train on that? Do I train and evaluate a separate policy per video?\\

        Partial-tile visibility &
        Any-overlap vs.\ overlap threshold vs.\ area-weighted coverage \\

        Parameters &
        Do I use the same as the paper: $P_{\max}$, $W$, $T_0$, $N_0$ as $20\,\mathrm{dBm}$, $10\,\mathrm{MHz}$, $1\,\mathrm{s}$, $-174\,\mathrm{dBm/Hz}$\\

        Critic count &
        One critic to start, or two (as in prior related work)?\\

        \bottomrule
    \end{tabular}
    \caption{Additional modeling decisions requiring confirmation.}
    \label{tab:decisions}
\end{table}
\noindent
\textbf{Constraint budgets.}
How should I define the constraint thresholds $\bar{D}$, $\bar{P}$, and $\bar{B}$ in the constrained MDP?

\[
\begin{aligned}
\bar{D} &: \text{ stall duration budget},\\
\bar{P} &: \text{power budget},\\
\bar{B} &: \text{number of enhanced tiles budget}.
\end{aligned}
\]

\noindent
Should these values be adopted from prior work,or selected experimentally or otherwise?
\section{Next Steps}

\begin{enumerate}
    \item Resolve the modeling questions listed above.
    \item Implement the Gymnasium environment using the formulated state,
    action, transition, constraint, and reward definitions.
    \item Train and validate a PPO baseline before introducing
    post-decision states and virtual experience.
\end{enumerate}

\end{document}
